% Part: first-order-logic
% Chapter: natural-deduction
% Section: provability-quantifiers

% Vérification des propriétés de dérivabilité utilisées pour les ensembles
% maximalement cohérents dans le chapitre sur la complétude.

\documentclass[../../../include/open-logic-section]{subfiles}

\begin{document}

\olfileid[fr]{fol}{ntd}{qpr}

\olsection{\usetoken{S}{derivability} et quantificateurs}

\begin{explain}
La démonstration du théorème de complétude exige aussi que les
règles de déduction naturelle donnent les propriétés de~$\Proves$
établies dans cette section.
\end{explain}

\begin{thm}
\ollabel{thm:strong-generalization} Si $c$ est une constante qui ne
figure ni dans $\Gamma$ ni dans $!A(x)$ et si $\Gamma \Proves !A(c)$,
alors $\Gamma \Proves \lforall[x][!A(x)]$.
\end{thm}

\begin{proof}
Soit $\delta$ une !!{derivation} de $!A(c)$ à partir de $\Gamma$.
En ajoutant une inférence \Intro{\lforall}, on obtient une
!!{derivation} de $\lforall[x][!A(x)]$. Puisque $c$ ne figure ni
dans $\Gamma$ ni dans $!A(x)$, la condition sur l'eigenvariable
est satisfaite.
\end{proof}

\begin{prop}
\ollabel{prop:provability-quantifiers}
Pour tout terme clos~$t$, on a :
\begin{tagenumerate}{prvEx,prvAll}
\tagitem{prvEx}{$!A(t) \Proves \lexists[x][!A(x)]$.}{}

\tagitem{prvAll}{$\lforall[x][!A(x)] \Proves !A(t)$.}{}
\end{tagenumerate}
\end{prop}

\begin{proof}
\begin{tagenumerate}{prvEx,prvAll}
\tagitem{prvEx}{Voici une !!{derivation}
  de~$\lexists[x][!A(x)]$ à partir de~$!A(t)$ :
\begin{prooftree}
\AxiomC{$!A(t)$}
\RightLabel{\Intro{\lexists}}
\UnaryInfC{$\lexists[x][!A(x)]$}
\end{prooftree}
}{}

\tagitem{prvAll}{Voici une !!{derivation} de~$!A(t)$
  à partir de~$\lforall[x][!A(x)]$ :
\begin{prooftree}
\AxiomC{$\lforall[x][!A(x)]$}
\RightLabel{\Elim{\lforall}}
\UnaryInfC{$!A(t)$}
\end{prooftree}
}{}
\end{tagenumerate}
\end{proof}

\begin{editorial}
La fermeture du terme~$t$, requise par les règles de ce système,
est rappelée explicitement dans l'énoncé de la proposition.
\end{editorial}

\end{document}
